\begin{pastexam}[label={px:gk3-110-ds-02}]{高考三級資訊處理 110 年 第二題}
（申論；25 分）

假設收銀機內銅板的集合 S=\{\$50, \$20, \$20, \$15, \$10, \$2, \$1, \$1, \$1\}，而預計找錢給顧客的金額
W=\$75。㈠請設計一個 Greedy（貪婪）的演算法，來解決找錢給顧客的問題，使得找給顧客金額 W 所使用的
銅板數量最少，並依此 Greedy 的演算法列出找給顧客金額 W=\$75 的過程。（15 分）㈡此 Greedy 演算法
適合使用何種資料結構來完成。（5 分）㈢此 Greedy 演算法的解法是否能保證為最佳解？請舉例說明。（5 分）
\end{pastexam}

\begin{pastsoln}
這題的銅板是\textbf{有限的一批}（例如 \$20 只有兩枚、\$2 只有一枚），和資結書第~19.5~節「每種面額
用不完」的找零錢稍有不同，貪心的想法一樣：每一步拿還放得下的最大銅板。

\textbf{㈠ 演算法}。 把收銀機裡的銅板由大到小排好，依序看每一枚：面額不超過還要找的金額就拿，
否則跳過。
\begin{pseudo}{有限銅板的貪心找零}
\Function{Greedy-Change}{$S, n, W$} \Comment{$S[1..n]$ 是收銀機裡每一枚銅板的面額}
    \State 把 $S$ 由大到小排序
    \State $\mathit{rest} \gets W$，$\mathit{used} \gets$ 空串列
    \For{$i \gets 1$ \textbf{to} $n$}
        \If{$S[i] \le \mathit{rest}$} \Comment{放得下就拿}
            \State 把 $S[i]$ 加進 $\mathit{used}$
            \State $\mathit{rest} \gets \mathit{rest} - S[i]$
        \EndIf
        \If{$\mathit{rest} = 0$}
            \State \Return $\mathit{used}$
        \EndIf
    \EndFor
    \State \Return \NIL \Comment{剩下的錢湊不出來}
\EndFunction
\end{pseudo}
排序 $O(n \log n)$，掃描 $O(n)$，共 $O(n \log n)$。

$W = 75$ 的過程（$S$ 已經由大到小）：
\begin{center}
\small
\begin{tabular}{cccl}
\toprule
步驟 & 這枚面額 & 還要找 & 動作 \\
\midrule
1 & 50 & 75 & $50 \le 75$，拿，剩 $25$ \\
2 & 20 & 25 & $20 \le 25$，拿，剩 $5$ \\
3 & 20 & 5 & $20 > 5$，跳過 \\
4 & 15 & 5 & $15 > 5$，跳過 \\
5 & 10 & 5 & $10 > 5$，跳過 \\
6 & 2 & 5 & $2 \le 5$，拿，剩 $3$ \\
7 & 1 & 3 & 拿，剩 $2$ \\
8 & 1 & 2 & 拿，剩 $1$ \\
9 & 1 & 1 & 拿，剩 $0$，結束 \\
\bottomrule
\end{tabular}
\end{center}
貪心找的是 $50 + 20 + 2 + 1 + 1 + 1 = 75$，共 $6$ 枚。

\textbf{㈡ 資料結構}。 貪心每一步都要「目前最大的那一枚」，所以適合：
\begin{itemize}
\item 由大到小排好的陣列（或串列），從頭掃到尾即可，上面的虛擬碼就是這樣做；
\item 或用最大堆積（優先佇列，資結書第~8~章）：建堆 $O(n)$，每次取出最大 $O(\log n)$。 銅板會
隨時放進收銀機、或只需要取出前幾枚時，堆積比整個重新排序方便。
\end{itemize}
如果面額種類少，也可以用「面額 $\to$ 枚數」的表（陣列依面額由大到小），一次算出某個面額要拿
$\min(\text{枚數}, \lfloor \mathit{rest} / \text{面額} \rfloor)$ 枚。

\textbf{㈢ 不保證最佳}。 這一題本身就是反例：
\[
75 = 50 + 15 + 10
\]
只要 $3$ 枚，貪心卻用了 $6$ 枚。 原因是貪心在第 2 步拿了 \$20 之後剩 \$5，而 \$5 只能用
\$2 和 \$1 湊；它在選 \$20 的當下看不到這個後果。 用不完的面額也有同樣的問題（資結書第~19.5~節）：
面額 $1, 5, 12$ 找 $15$，貪心拿 $12 + 1 + 1 + 1$ 共 $4$ 枚，$5 + 5 + 5$ 只要 $3$ 枚。 只有像新臺幣
$1, 5, 10, 50$ 這種「標準的」面額組合，而且每種面額都夠用時，貪心才一定最少。 一般情況要用動態
規劃（有限枚數時是 0/1 背包式的 DP，資結書第~18.5~節）。
\end{pastsoln}
